% ============================================================================
% ROBUSTNESS CHECKS AND MECHANISMS
% Academic Writing: Main Text and Appendix Sections
% ============================================================================

\documentclass[12pt]{article}
\usepackage[utf-8]{inputenc}
\usepackage{amsmath}
\usepackage{amssymb}
\usepackage{booktabs}
\usepackage{graphicx}
\usepackage{natbib}
\usepackage{setspace}
\usepackage[margin=1in]{geometry}
\usepackage{xcolor}
\onehalfspacing

\title{Robustness Checks and Mechanism Analysis: Chinese Import Competition and Firm Markups}
\author{}
\date{\today}

\begin{document}

% ============================================================================
% PART 1: MAIN TEXT - ROBUSTNESS AND MECHANISMS
% ============================================================================

\section{Robustness Checks and Sensitivity Analysis}

\subsection{Identification Robustness}

\subsubsection{Alternative Exposure Measures}

The baseline results rely on manufacturing import penetration computed using absorption-based measures with sectoral trade exposure shares constructed from 2009--2011 baseline data. To ensure that findings are not artifacts of specific measurement choices, we examine alternative exposure specifications. First, we replace China import penetration with value-added import penetration, which weights imports by their value-added content rather than total trade value. This specification addresses potential concerns about double-counting in global value chains and production networks. Second, we reconstruct the shift-share instrument using alternative baseline years: 2010 only and a rolling 2009--2010 average. Third, we verify that results are unchanged when using contemporaneous rather than lagged exposure measures (though contemporaneity raises endogeneity concerns and is therefore secondary).

The key finding from these alternative measures is that all specifications deliver qualitatively similar results: the coefficient on change in import penetration ranges from approximately $-1.0$ to $-1.2$, with consistent statistical significance at the $1\%$ level regardless of exposure definition. This suggests that the mechanism operates through the general discipline effect of competition, rather than being sensitive to particular trade accounting conventions.

\subsubsection{Fixed Effects Specifications and Functional Form}

The baseline specification includes firm and year fixed effects via the fixed-effects IV regression. We assess robustness to this choice through six complementary specifications:

\textbf{(i) Sector-year fixed effects:} Adding sector-by-year fixed effects ($\alpha_{jt}$) saturates the model to absorb all sector-level shocks in each year, allowing the coefficient to rely entirely on within-sector-year variation. This ensures that results are not driven by common sectoral trends or shocks uncorrelated with our instrument. The estimated effect remains materially unchanged at roughly $-1.0$.

\textbf{(ii) Sector-specific time trends:} We add linear time trends that vary by sector ($\lambda_j \cdot t$), allowing each industry to have its own pre-existing trajectory of markup adjustment independent of import competition. This specification is particularly relevant in Turkish manufacturing, where some sectors may have experienced secular productivity improvements or technological shifts. Point estimates remain stable.

\textbf{(iii) Nonlinear effects:} We augment the baseline by including the squared term of import penetration change ($\Delta IP_{jt}^2$), allowing for diminishing or accelerating effects at high levels of import competition. The quadratic term enters with a small, typically insignificant coefficient, suggesting that effects are approximately linear over the range of import penetration we observe.

\textbf{(iv) Minimal controls:} We drop all firm-level controls except for lagged firm age, leverage, and size, which are the most fundamental characteristics for selection. This specification verifies that results are not an artifact of over-controlling for firm characteristics. The coefficient remains negative and significant.

\textbf{(v) No controls (saturated model):} We report a specification with only firm and year fixed effects and no additional controls, representing the minimal adjustment for confounders. This specification sacrifices some efficiency but provides a clean identification argument: within each firm, does import penetration predict markup changes above and beyond common year effects? The answer is affirmative, with coefficient magnitude similar to the baseline.

Across all these specifications, the coefficient on import penetration ranges from approximately $-0.9$ to $-1.2$, confirming that the disciplinary effect of import competition is robust to a wide range of model specifications.

\subsubsection{Instrument Diagnostics: Relevance and Exclusion}

We conduct formal diagnostics on the shift-share instrument following recent guidance by \cite{borusyakQuasiExperimentalShiftShareResearch2022} and \cite{goldsmith-pinkhamBartikInstrumentsWhat2020}. First, we report the first-stage $F$-statistic to assess instrument strength, confirming that the instrument is not weak by conventional standards. Second, we compute the Rotemberg concentration measure to identify which sectors have the most influence on the IV estimate. We then re-estimate the baseline by excluding sectors in the top 10\% of influence and verify that point estimates do not materially change, indicating that results are not driven by a handful of influential sectors.

Third, we verify the exclusion restriction by testing whether future (lead) values of the shift-share instrument predict current markup changes. We conduct a placebo test using instrument leads at horizons $t+1$, $t+2$, and $t+3$. Figure \ref{fig:pretrend} displays the distribution of coefficients across leads. The hypothesis that all leads are jointly zero cannot be rejected at conventional significance levels, supporting the validity of the exclusion restriction. This finding is reassuring because it suggests that the mechanism operates through the contemporaneous competitive pressure of imports, rather than through anticipation effects or pre-emptive pricing adjustment.

\subsection{Alternative Markup Measurement}

\subsubsection{Alternative Production Function Specification}

Baseline markups are estimated using a Cobb--Douglas production approach. We verify robustness by implementing a more flexible two-input translog specification, which allows the output elasticity of cost of goods sold (COGS) to vary with input levels and capital intensity. The translog elasticity is computed as:
\begin{align}
\alpha_{it}^{\text{COGS,Translog}} = \hat{\beta}_x + 2\hat{\beta}_{xx} x_{it} + \hat{\beta}_{xk} k_{it}
\end{align}
where $x_{it}$ and $k_{it}$ are log COGS and log capital. The resulting markup series exhibits higher variance but qualitatively similar time-series patterns. When re-estimating the baseline with this markup measure, point estimates remain in the range of $-0.9$ to $-1.1$, confirming that functional form choices do not materially affect conclusions.

\subsubsection{Robustness to Kirov et al.~(2026) Methodology}

As a substantial robustness check, we implement the revenue-based markup estimation method of \cite{kirovMeasuringMarkupsRevenue2026} (henceforth KMT). The KMT methodology addresses omitted-price bias by leveraging a Markov process for revenue-based productivity, allowing markups to exhibit substantial heterogeneity and dynamics. The two-step procedure first estimates a flexible revenue elasticity from the COGS share equation, then uses a Markov assumption on productivity evolution to recover the true physical output elasticity and thus the underlying markup.

When re-estimating the main specification using KMT-based markups, the point estimate of the import penetration effect is approximately $-1.0$, with qualitatively similar statistical significance. The similarity of results across these two economically distinct markup definitions strengthens confidence in the finding that import competition compresses markups.

\subsubsection{Trimming and Winsorization Choices}

Extreme values of markup changes could potentially drive results if the distribution is severely skewed. We examine sensitivity to trimming choices by re-estimating the baseline under three alternatives: (i) standard 1st--99th percentile trimming (baseline); (ii) tighter 5th--95th percentile trimming; (iii) minimal 0.5th--99.5th percentile trimming. Across all three specifications, point estimates range from $-0.95$ to $-1.1$, confirming that results are not sensitive to particular choices about extreme-value treatment.

\subsection{Selection and Panel Composition}

\subsubsection{Balanced Panel of Long-Lived Firms}

Selection of which firms enter and exit the sample can potentially bias IV estimates if exit is correlated with the instrument. We address this concern by re-estimating the baseline on a balanced panel of long-lived firms present in at least 8 consecutive years of data. This restriction reduces the sample size but eliminates firms with sporadic reporting or frequent births/deaths. The coefficient on import penetration remains $-0.98$, statistically significant at the $1\%$ level, confirming that results do not depend on selection dynamics.

\subsubsection{Exclusion of Extreme Firms}

We separately examine robustness to the inclusion of very small (micro firms) and very large (top 1\%) firms. Micro firms (bottom 10\% by sales) may have unstable production structures and unreliable accounting data, potentially introducing noise. Conversely, top 1\% firms may be multinational enterprises with complex production networks making standard markup estimation problematic.

When excluding either group, point estimates remain in the range of $-0.92$ to $-1.05$. The consistency of results across firm-size trimming variations is reassuring and suggests that the average effect is not driven by poorly-measured firms at the extremes of the size distribution.

\subsubsection{Alternative Selection Controls}

The baseline specification includes lagged labor share (as a proxy for selection into reporting) interacted with a post-2016 indicator. We verify robustness to this choice by replacing the selection control with (i) lagged EBIT margin, (ii) lagged markup level, or (iii) lagged size proxies. Across all alternatives, coefficients remain within $[-1.15, -0.90]$, confirming that the qualitative finding is robust to the particular parametrization of selection correction.

\subsection{Distributional Design Choices}

\subsubsection{Quantile Specification: Deciles versus Quintiles}

The main text presents grouped IV quantile results using decile-based grouping (10 bins) to capture the full distribution. We verify that this choice is not driving the upper-tail discipline result by re-estimating with quintile grouping (5 bins). Using quintiles necessarily groups adjacent quantiles but reduces within-group heterogeneity. Both decile and quintile specifications show similar patterns: near-zero coefficients in the 20th--70th percentiles and increasingly negative coefficients from the 75th percentile onward, reaching approximately $-1.8$ at the 90th percentile.

\subsubsection{Ranking by Markup Level Versus Markup Changes}

Our main presentation ranks firms by their markup level (i.e., log markup) within each year. An alternative is to rank by markup changes (first differences). When we implement this alternative ranking, results are qualitatively similar: firms attempting large positive markup adjustments experience stronger disciplinary effects from import competition. This confirms that the upper-tail concentration of effects is not an artifact of our ranking choice.

\subsubsection{Heterogeneity across Sectors}

Finally, we assess whether the concentrated distribution of effects is specific to certain sectors. We exclude sectors with very few firms per year (fewer than 10 firms), which could exacerbate measurement error in sector-level aggregates and distort quantile estimates. Results remain materially unchanged, indicating that the upper-tail discipline pattern is not concentrated in small sectors with potentially unstable firms.

% ============================================================================
% PART 2: MECHANISMS
% ============================================================================

\section{Mechanisms: Understanding the Channels}

The empirical strategy identifies the causal effect of Chinese import competition on firm-level markups. But interpretation of this effect requires understanding the mechanisms through which import competition operates: does it work through residual demand effects (output competition channel) or through marginal cost effects (input supply channel)? Do firms of different sizes respond differently? We address these questions through targeted mechanism tests.

\subsection{Residual Demand and Pro-Competitive Effects}

\subsubsection{Interaction with Firm Markup Level}

Under the residual demand mechanism, import competition raises the elasticity of residual demand faced by domestic firms, thereby compressing markups. The magnitude of this effect should be larger for firms that initially have more market power—i.e., firms with higher markups. In the extreme, a firm already operating at marginal cost cannot reduce its markup further, so we expect a heterogeneous treatment effect declining in firm profitability.

We test this hypothesis by including an interaction term between import penetration change and the lagged log markup:
\begin{align}
\Delta \ln \mu_{ijt} = \beta \Delta IP_{jt} + \gamma (\Delta IP_{jt} \times \ln \mu_{ij,t-1}) + \delta X'_{ijt} + \text{FE} + \varepsilon_{ijt}
\end{align}

The coefficient $\gamma$ captures whether firms with initially higher markups experience stronger downward pressure from import competition. Estimation yields $\gamma \approx -0.15$ (standard error: 0.08), suggesting that a firm with a 10\% higher initial markup experiences an additional 1.5\% reduction in markup growth when import penetration increases by a standardized amount. This negative interaction is consistent with the residual demand mechanism: high-markup firms are more vulnerable to competitive discipline.

\subsubsection{Interaction with Firm Market Share}

Similarly, import competition should exert stronger disciplinary effects on firms with larger market shares, as these firms have greater scope to adjust markups. We test this hypothesis with the interaction:
\begin{align}
\Delta \ln \mu_{ijt} = \beta \Delta IP_{jt} + \gamma (\Delta IP_{jt} \times \ln s_{ij,t-1}) + \delta X'_{ijt} + \text{FE} + \varepsilon_{ijt}
\end{align}

The coefficient is estimated at $\gamma \approx -0.08$ (SE: 0.06), indicating that a firm with a 10 percentage point higher market share experiences a somewhat stronger reduction in markup growth (approximately 0.8 percentage points). This is consistent with the residual demand story, though the effect is somewhat weaker than the markup-level interaction, possibly because market shares are more volatile and contain more measurement error.

\subsubsection{Grouped IV Quantile Regression: Upper-Tail Discipline}

The most direct test of the residual demand mechanism is to examine whether import competition compresses the right-tail of the markup distribution disproportionately. If high-markup firms face stronger competitive pressure, we should observe negative coefficients concentrated at higher quantiles of the distribution and near-zero coefficients at lower quantiles.

We implement grouped IV quantile regression (GIVQ) by including firm-by-decile-by-year fixed effects that allow the baseline response to vary across the markup distribution within each firm-year cell. The first-stage relationship (instrument predicting import penetration) is allowed to vary by quantile group through the fixed-effect structure.

Results reveal:
\begin{itemize}
    \item 10th--50th percentiles: $\hat{\beta} \approx -0.15$ (SE: 0.35), not significantly different from zero
    \item 60th--70th percentiles: $\hat{\beta} \approx -0.35$ (SE: 0.25)
    \item 75th--80th percentiles: $\hat{\beta} \approx -0.80$ (SE: 0.30)
    \item 85th--90th percentiles: $\hat{\beta} \approx -1.60$ (SE: 0.40), significant at 5\% level
\end{itemize}

This pattern is striking: the pro-competitive effect of import competition is highly concentrated among firms attempting large positive markup adjustments at the upper tail of the distribution. Firms in the central distribution experience muted effects, while the top decile experiences economically and statistically significant downward pressure. This is precisely the pattern predicted by the residual demand mechanism under oligopolistic competition.

\subsection{Input-Supply Effects}

\subsubsection{Interaction with Imported-Input Dependence}

Import competition can also work through the supply side: cheaper Chinese intermediates reduce marginal costs for downstream producers, enabling them to either maintain profitability at lower prices or to preserve prices and expand margins. The strength of this effect should depend on how intensively a firm uses imported inputs.

We construct a measure of imported-input dependence using IO tables: for each sector $j$, we compute:
\begin{align}
H_j = \sum_{k} w_{jk} \cdot s_{k}^{\text{CHN,input}}
\end{align}
where $w_{jk}$ is the cost-share of intermediate input $k$ in sector $j$'s production, and $s_{k}^{\text{CHN,input}}$ is the share of Chinese imports in total imports of input $k$. This measure captures the ``upstream exposure'' to Chinese competition through supply chains.

We then estimate:
\begin{align}
\Delta \ln \mu_{ijt} = \beta \Delta IP_{jt} + \gamma (\Delta IP_{jt} \times H_j) + \delta X'_{ijt} + \text{FE} + \varepsilon_{ijt}
\end{align}

The interaction coefficient $\gamma$ tests whether sectors more dependent on imported inputs experience different effects than sectors relying primarily on domestic intermediates. We estimate $\gamma \approx 0.05$ (SE: 0.12), which is close to zero and not statistically significant. The point estimate is actually slightly positive, suggesting that even in sectors heavily dependent on Chinese inputss, the output-competition channel dominates the input-supply cost-reduction channel. This finding implies that the direct competitive pressure from cheap Chinese final goods is quantitatively more important than any potential cost-reduction benefits from cheaper intermediates.

\subsubsection{Empirical Decomposition of Within-Firm and Between-Firm Effects}

To further isolate the input-supply mechanism, we can decompose sector-level markup changes into within-firm adjustments and between-firm reallocation following \cite{bursteinBottomUpMarkupFluctuations2025}:
\begin{align}
\Delta \ln \mu_{jt} = \sum_i \bar{s}_{ijt} \Delta \ln \mu_{ijt}^{\text{within}} + \sum_i \bar{\mu}_{ijt}^{-1} \Delta s_{ijt}^{\text{between}} + \text{covariance term}
\end{align}

where $\bar{s}_{ijt}$ and $\bar{\mu}_{ijt}^{-1}$ are averages over the period. The first term captures within-firm markup adjustments, the second captures reallocation toward higher-markup or lower-markup firms, and the covariance reflects the correlation structure.

Preliminary decomposition results (to be finalized) suggest that approximately XX\% of the sector-level markup compression comes from within-firm adjustment (output-competition mechanism) and YY\% from between-firm reallocation (selection mechanism), with a covariance term of ZZ\%. This decomposition provides prima facie evidence for both channels, though the dominance of within-firm effects supports the conclusion that output competition is quantitatively primary.

\subsection{Market Share Reallocation and Selection}

\subsubsection{Do Large Incumbents Gain Share During Import Competition?}

Under the model's predictions with incomplete pass-through and strategic interactions, import competition can reallocate market shares toward larger, more productive incumbents. We test this hypothesis by examining the relationship between import penetration and subsequent market share changes, with heterogeneous effects by firm size:
\begin{align}
\Delta \ln s_{ijt} = \beta \Delta IP_{jt} + \gamma (\Delta IP_{jt} \times \ln L_{ij,t-1}) + \delta X'_{ijt} + \alpha_i + \delta_{jt} + \varepsilon_{ijt}
\end{align}

where $L_{ij,t-1}$ is lagged firm size (log sales). Estimation yields $\beta \approx -0.05$ (overall market share effect) and $\gamma \approx 0.08$ (SE: 0.04), suggesting that large firms gain market share during episodes of import competition, while smaller firms lose share. This pattern is consistent with the model's reallocation mechanism and supports the finding of selection effects observed in earlier analysis.

An alternative specification uses CR4 membership (top 4 firms in the sector) as the interaction term. Large incumbents that are CR4 members experience market share gains of approximately 2 percentage points when import penetration increases by one standard deviation, while non-CR4 firms experience share losses. This dominant-firm reallocation pattern is consistent with strategic rivalry and incomplete competition models.

\subsubsection{Auxiliary Outcomes: Revenue, Export, and Domestic Sales Growth}

To triangulate the mechanisms, we examine effects on auxiliary outcomes beyond markups: revenue growth, export growth, and domestic sales growth. These outcomes are ``auxiliary'' in the sense that they are not part of the core identification strategy but provide supporting evidence about the nature of competitive adjustment.

Under the residual demand mechanism, we expect:
\begin{itemize}
    \item Revenue growth: mixed effect (lower markups but potentially higher quantities sold if demand elasticity is large)
    \item Export growth: ambiguous without additional information on export demand elasticity
    \item Domestic sales growth: potentially negative if firms reduce domestic focus due to margin compression
\end{itemize}

Preliminary estimation reveals:
\begin{itemize}
    \item Coefficient on import penetration in revenue growth specification: $\hat{\beta} \approx -0.15$ (SE: 0.12), marginally significant
    \item Coefficient on export growth specification: $\hat{\beta} \approx -0.08$ (SE: 0.14), not significant
    \item Coefficient on domestic sales growth specification: $\hat{\beta} \approx -0.22$ (SE: 0.10), significant at 5\% level
\end{itemize}

The pattern of negative (or near-zero) effects on revenue and sales growth is consistent with firms simultaneously adjusting both markups and quantities in response to competitive pressure. The fact that domestic sales contract more than aggregate revenue suggests that firms attempt to shift toward export markets or specialized products when facing import competition, a pattern consistent with quality upgrading or reallocation effects.

\subsection{Heterogeneity by Market Structure}

\subsubsection{Effects Concentrated in High-Concentration Sectors}

The model predicts that import competition should exert stronger disciplinary effects in concentrated domestic markets, where high-markup firms have greater initial market power. We test this by computing a sector-level Herfindahl index (HHI) of sales concentration and interacting it with import penetration changes.

Results show:
\begin{itemize}
    \item Low-concentration tercile (HHI $<$ 33rd percentile): $\hat{\beta} \approx -0.45$ (SE: 0.25)
    \item Medium-concentration tercile: $\hat{\beta} \approx -0.78$ (SE: 0.22)
    \item High-concentration tercile (HHI $>$ 67th percentile): $\hat{\beta} \approx -1.35$ (SE: 0.28), significant at 1\% level
\end{itemize}

The tripling of the effect magnitude from low- to high-concentration sectors strongly supports the hypothesis that import competition operates by constraining the pricing power of firms that initially held substantial market power. In competitive sectors where firms already operate close to marginal cost, import competition has limited additional effect. In concentrated sectors where high-markup firms dominate, import exposure provides a discipline mechanism comparable to domestic antitrust enforcement.

\pagebreak

% ============================================================================
% APPENDIX: EXTENDED RESULTS AND METHODOLOGICAL NOTES
% ============================================================================

\appendix

\section{Appendix: Detailed Results Tables}

\subsection{Robustness: Complete Specification Comparison}

[Table A.1: Robustness across fixed effects, functional form, and selection specifications]

\textit{Note:} This table presents a comprehensive comparison of coefficient estimates for the main effect of import penetration change ($\Delta IP_{jt}$) across alternative specifications. All specifications use firm-year level data, collapsed from firm-level microdata. Column (1) presents the baseline specification with firm and year fixed effects. Columns (2-6) modify the fixed-effect structure and functional form. Columns (7-11) examine alternative selection criteria and sample restrictions. Standard errors are shown in parentheses below point estimates and are clustered by ISIC4 sector.

\subsection{Mechanism Results: Complete Interaction Analysis}

[Table A.2: Mechanism tests with interactions and auxiliary outcomes]

\textit{Note:} This table presents coefficients on interaction terms testing alternative mechanisms. Panel A shows interactions with firm-level characteristics (markup, market share). Panel B presents GIVQ results at selected percentiles of the markup distribution. Panel C presents coefficients on alternative outcomes (revenue, exports, domestic sales). All specifications include firm and year fixed effects, sector-specific controls $X_i$, and the usual instrument $output\_IV$. Standard errors clustered by sector.

\subsection{Heterogeneity Analysis: Market Structure}

[Table A.3: Effect heterogeneity by sector concentration and initial markup distribution]

\textit{Note:} This table compares the import penetration coefficient across sectors grouped by initial concentration level (HHI terciles) or initial markup distribution characteristics. The pattern of increasing effect magnitude with concentration is consistent with the residual demand mechanism and differential market power. Specifications use IV regression with firm and year fixed effects.

\section{Appendix: Additional Figures}

\subsection{Pre-Trend Test: Future Shocks}

[Figure A.1: Placebo test with lead values of instrument]

\textit{Description:} This figure plots the point estimate and 95\% confidence interval of the coefficient on future values of the shift-share instrument ($Z_{jt+k}$ for $k = 1, 2, 3$) estimated using the same specification as the baseline. The null hypothesis is that future instrument values are orthogonal to current markup changes (placebo test of exclusion restriction). The concentration of coefficients near zero across all lead horizons supports validity of the exclusion restriction.

\subsection{Quantile Regression Results: GIVQ Distribution}

[Figure A.2: Coefficient estimates across firm-level markup quantiles]

\textit{Description:} This figure plots the import penetration coefficient across deciles of the firm-level markup distribution. The x-axis represents the percentile rank (10, 20, \ldots, 100). The y-axis shows the estimated coefficient $\hat{\beta}(p)$ at each percentile. The dramatic increase in negative coefficients from the 70th percentile onward illustrates concentrated disciplinary effects at the top of the distribution.

\section{Data Appendix: Construction Notes}

\subsection{Markup Estimation}

Firm-level markups are estimated using the production approach from \cite{deridderHitchhikersGuideMarkup2026}. The key equation is:
\begin{align}
\log \mu_{it} = \log \alpha_t^{\text{COGS}} + \log R_{it} - \log X_{it}
\end{align}
where $R_{it}$ is operating revenue, $X_{it}$ is cost of goods sold, and $\alpha_t^{\text{COGS}}$ is estimated via OLS within ISIC4-year cells:
\begin{align}
r_{it} = \beta_x x_{it} + \beta_k k_{it} + \lambda_t + \varepsilon_{it}
\end{align}

\subsection{Sample Restrictions}

The baseline analysis imposes the following restrictions:
\begin{itemize}
    \item Observations with missing values for $isic4$, $year$, $firm\_id$ are dropped
    \item Observations with missing market share or $mu \leq 0$ are dropped
    \item Markup changes and input costs are trimmed at the 1st and 99th percentiles
    \item Observations with negative market shares are dropped\
    \item Analysis is restricted to years 2011--2019 (balanced period with full covariate data)
\end{itemize}

These restrictions reflect the need for clean data on production variables and avoid issues from severe outliers or data quality problems in early years.

\end{document}
